\section{问题重述}

\subsection{任务背景与统一约束}
题目给出三枚来袭导弹、五架无人机及一个圆柱形真实目标。导弹均以 $300\,\mathrm{m/s}$ 的恒定速率飞向位于坐标原点的假目标；无人机受领任务后可瞬时确定水平航向和飞行速度，此后保持高度并作匀速直线运动，速度范围为 $70\sim140\,\mathrm{m/s}$。无人机投放的烟幕弹在重力作用下运动，起爆后形成有效半径 $10\,\mathrm m$ 的球状烟幕云团，云团在 $20\,\mathrm s$ 有效期内以 $3\,\mathrm{m/s}$ 匀速下沉。同一无人机相邻两次投弹至少间隔 $1\,\mathrm s$。上述条件均直接来自题面\cite{ref1}。

真实目标不是空间点，而是半径 $7\,\mathrm m$、高度 $10\,\mathrm m$ 的有限圆柱。因此，各问的共同前提是：烟幕是否真正阻断导弹到\emph{整个有限目标可见部分}的视线，而不是仅遮住目标中心或与某条无限直线相交。

\subsection{五个问题的任务链}
题目五问沿“固定策略评价--单弹优化--同机多弹--多机协同--多导弹联合”的方向逐级扩展：
\begin{enumerate}[label=\arabic*.]
  \item 在 FY1 的飞行方向、速度、投弹时刻和起爆延迟均给定时，计算单枚烟幕弹对 M1 的有效遮蔽时长；
  \item 释放 FY1 的航向、速度、投弹时刻和起爆延迟，求单机单弹的最优投放策略；
  \item FY1 连续投放三枚烟幕弹，在共享航向、共享速度及相邻投弹间隔约束下设计三弹时序；
  \item FY1--FY3 各投放一枚烟幕弹，联合确定三架无人机的航迹与起爆参数，使对 M1 的有效遮蔽尽可能延长；
  \item 五架无人机各至多投放三枚烟幕弹，对 M1--M3 实施联合干扰，并给出完整可执行的投放方案及三枚导弹的有效遮蔽结果。
\end{enumerate}

\section{问题分析}

五个问题共用同一组运动规律和目标几何，但数学难点逐问增加。全题首先需要建立一个能在有限圆柱目标上统一使用的\textbf{完整遮蔽评价器}；随后只改变决策变量的组织方式和协同层级。由此可将全题组织为“公共运动关系--有限目标几何判定--单烟幕参数优化--多烟幕/多平台协同--多导弹联合评价”的递进链。

\subsection{问题一：从有限目标构造完整遮蔽判据}
本问没有优化变量，关键不是运动轨迹本身，而是“有效遮蔽”的几何定义。对任一时刻，导弹到目标可见点形成一束\emph{有限视线段}；单枚烟幕只有同时截获这束视线中的全部成员，才能形成完整遮蔽。因此需要先定义点到有限线段的距离，再对有限圆柱的可见视线集合取最不利值。该评价器一旦建立，可直接作为后续四问的统一几何核。

\subsection{问题二：四维参数化下的单机单弹优化}
FY1 一旦确定航向和速度就不再调整，投放点与起爆点也分别由投弹时刻和起爆延迟唯一决定，因此无需把空间坐标重复作为独立决策变量。问题二可压缩为“航向--速度--投弹时刻--起爆延迟”四个连续变量。其目标值由完整遮蔽时间集合的长度决定，随参数变化会出现区间生成、消失和边界切换，因而属于非凸、非光滑的连续优化问题。

\subsection{问题三：共享平台约束下的三弹联合遮蔽}
三枚烟幕共享同一架无人机的航向和速度，同时受投弹间隔约束。多烟幕协同不能简单把三枚烟幕的独立遮蔽时长相加：在同一时刻，不同烟幕可以分别遮挡不同的目标视线，因此应对每条视线先选择最有利的有效烟幕，再对全部可见视线取最坏情形。该结构把“同机时序约束”和“空间互补遮蔽”耦合为一个八维连续问题。

\subsection{问题四：不同初始位置下的三机时空协同}
三架无人机的航向和速度互不共享，但初始位置差异显著。与问题三相比，本问解除同机三弹的共享平台约束，却增加三组独立航迹变量。需要判断三架无人机的有效区间是依靠同时空间拼接扩展，还是在不同时间段形成接力。因而除了求最大遮蔽时长，还应分析各机独立区间与联合区间之间的集合关系。

\subsection{问题五：五机多弹对三导弹的联合决策}
问题五同时包含五组固定航迹、最多十五个投弹事件和三枚来袭导弹。由于任一有效烟幕都可能依据实时几何关系帮助任一导弹，若预先设置排他的“无人机--导弹”分配变量，反而可能截断真实协同。因此正式模型应保留全部烟幕对三枚导弹的自由几何贡献，并以三枚导弹各自有效时间的总量作为主效能指标；同时另行统计墙钟并集，以解释双重、三重同时干扰产生的协同增益。

综上，全题技术路线如图\ref{fig:route}所示。图中公共几何判据从问题一向后继问题持续传递，而问题三、问题四分别提供“同机多弹”和“多机协同”的结构基础，最终汇入问题五。

\begin{figure}[tbp]
  \centering
  \includegraphics[width=0.95\textwidth]{F00_全题技术路线图.pdf}
  \caption{全题问题链与技术路线}
  \label{fig:route}
\end{figure}

\section{模型假设}

题面已明确的导弹速度、无人机匀速等高直线飞行、烟幕有效半径与下沉速度属于已知事实，不再重复列为模型假设。为补足题面未规定但计算所必需的物理细节，作如下简化：
\begin{enumerate}[label=\arabic*.]
  \item \textbf{烟幕弹自由飞行阶段忽略空气阻力与风场。} 烟幕弹脱离无人机瞬间继承载机的水平速度，竖直初速度取零；该假设使投放到起爆阶段可用经典抛体运动描述。若存在强风或显著气动阻力，起爆点会产生额外水平/竖直偏移。
  \item \textbf{起爆后的有效云团采用硬边界球模型。} 在题给 $20\,\mathrm s$ 有效期内，以云团中心 $10\,\mathrm m$ 内为有效遮蔽区域；除题给 $3\,\mathrm{m/s}$ 竖直下沉外，不再加入水平漂移、半径扩张和浓度梯度。该简化与题面给出的有效范围直接对应，但不能描述真实烟幕的非均匀扩散。
  \item \textbf{多枚烟幕之间不存在负向相互作用。} 不考虑云团碰撞、稀释或遮蔽能力相互削弱；多烟幕仅通过几何覆盖形成联合效应。因而增加一枚可行烟幕不会使已有视线的几何遮蔽能力变差。
  \item \textbf{飞行与起爆过程按确定性计划执行。} 不考虑导航、投放与引信控制误差，题面坐标和时刻视为确定量。该假设用于求题设条件下的基准最优策略，实际部署时可在本文模型外进一步加入执行误差与鲁棒裕度。
\end{enumerate}

\section{主要符号说明}
\begingroup
\normalsize
\renewcommand{\arraystretch}{1.22}
\setlength{\tabcolsep}{6pt}
\begin{center}
\begin{tabularx}{\textwidth}{>{\centering\arraybackslash}p{0.17\textwidth} X >{\centering\arraybackslash}p{0.18\textwidth}}
\toprule
\textbf{符号} & \textbf{含义} & \textbf{单位} \\
\midrule
$t$ & 连续时间 & s \\
$\vec m_j(t)$ & 导弹 $M_j$ 在时刻 $t$ 的位置向量 & m \\
$\vec f_i(t)$ & 无人机 FY$i$ 在时刻 $t$ 的位置向量 & m \\
$\theta_i$ & FY$i$ 的水平航向角 & rad 或 $^\circ$ \\
$v_i$ & FY$i$ 的飞行速度 & m/s \\
$t_{ik}$ & FY$i$ 第 $k$ 枚烟幕弹的投放时刻 & s \\
$\delta_{ik}$ & FY$i$ 第 $k$ 枚烟幕弹的起爆延迟 & s \\
$\tau_{ik}$ & FY$i$ 第 $k$ 枚烟幕弹的起爆时刻 & s \\
$\vec c_{ik}(t)$ & 第 $(i,k)$ 枚有效烟幕云团中心位置 & m \\
$\mathcal T$ & 真实目标对应的有限闭圆柱集合 & --- \\
$\mathcal V_j(t)$ & 从 $M_j$ 观察到的目标可见点集合 & --- \\
$d_{\rm seg}$ & 点到有限视线段的最短欧氏距离 & m \\
$D_j(t)$ & 对 $M_j$ 的完整遮蔽最坏距离 & m \\
$E_j$ & $M_j$ 的有效遮蔽时间集合 & --- \\
$H_j$ & $M_j$ 的有效遮蔽时长 $\mu(E_j)$ & s \\
$H_\Sigma$ & 三枚导弹有效遮蔽时长之和 & 导弹$\cdot$s \\
$H_\cup$ & 三枚导弹有效时间集合的墙钟并集长度 & s \\
\bottomrule
\end{tabularx}
\vspace{0.12em}
\begin{minipage}{\textwidth}
{\small\color{cumcmgray}\raggedright 注：仅列出跨问反复出现且评委需要回查的主要符号。局部投影参数、算法迭代索引和临时变量均在首次出现处定义。$H_j$ 始终按导弹编号索引，不再用 $H_1,H_2,H_3,H_4$ 表示问题编号结果。\par}
\end{minipage}
\end{center}
\endgroup
